\section{防御机制}

\subsection{防御原则}

\begin{importantbox}{三大防御原则}
\begin{enumerate}
    \item \textbf{纵深防御}（Defense in Depth）：瑞士奶酪模型——多层防御，即使某层失败，其他层仍可阻止攻击。层次包括：输入清洗 $\to$ 模型加固 $\to$ 动作策略执行 $\to$ 监控与异常检测
    \item \textbf{最小权限与权限分离}（Least Privilege \& Privilege Separation）：系统只应拥有完成其功能所需的最小权限
    \item \textbf{安全设计}（Safe/Secure by Design）：理想情况下通过形式化验证提供可证明的安全保证
\end{enumerate}
\end{importantbox}

\subsection{具体防御机制}

\subsubsection{模型加固}

在模型开发的各个阶段进行加固：
\begin{itemize}
    \item 数据清洗和数据准备阶段
    \item 安全预训练阶段
    \item 后训练对齐阶段（Post-training Alignment）
    \item 机器反学习（Machine Unlearning）
\end{itemize}

\subsubsection{输入清洗护栏}

对 LLM 的输入进行检查：匹配预定义标准、转义特殊字符、规范化为标准格式，以阻断恶意输入。

\subsubsection{动作策略执行}

\begin{importantbox}{可编程权限控制框架}
Dawn Song 团队开发了 LLM Agent 的可编程权限控制框架：
\begin{itemize}
    \item 对 Agent 的每个工具调用进行策略检查
    \item 例如：禁止 \texttt{delete\_db} 操作，仅允许加载特定数据库
    \item 即使 RAG 知识库被投毒导致 Agent 尝试删除数据库，策略执行层也会阻断该操作
\end{itemize}
\end{importantbox}

银行应用案例展示了如何通过策略限制转账金额上限和收款人范围，防止 Agent 被操纵执行超出授权的金融操作。

\subsubsection{输出护栏}

对 LLM 输出进行检查，阻止有害、有毒或不当内容到达用户或外部系统。

\subsubsection{监控与异常检测}

持续监控 Agent 行为，检测异常模式，及时发现并响应潜在攻击。

\subsection{本章小结}

Agentic AI 的防御需要多层次、多机制的综合策略，从模型加固到输入清洗、动作控制、输出过滤和持续监控，形成完整的纵深防御体系。

